\section{Tổng kết và Hướng phát triển}

\subsection{Kết quả đạt được}
Đồ án đã thiết kế và triển khai thành công Hệ thống Quản lý Thư viện Trực tuyến tích hợp Trí tuệ Nhân tạo. Những kết quả chính bao gồm:
\begin{itemize}
    \item \textbf{Phát triển nền tảng Full-Stack hiện đại:} Xây dựng hệ thống ổn định với Frontend (React.js, Vite), Backend (Node.js, Express, Prisma, PostgreSQL) và AI Microservice (Python, FastAPI). Triển khai thành công trên môi trường đám mây (Vercel, Google Cloud Run).
    \item \textbf{Tích hợp tính năng AI hiệu quả:} Cài đặt thành công tìm kiếm ngữ nghĩa (Semantic Search) kết hợp lai (Hybrid) với từ khóa, Trợ lý ảo AI (Chatbot) ứng dụng RAG để giải đáp nội quy thư viện, hệ thống gợi ý sách cá nhân hóa, và tự động điền thông tin sách qua ISBN.
    \item \textbf{Số hóa toàn diện nghiệp vụ thư viện:} Hoàn thiện các quy trình quản lý danh mục sách, kho ấn bản vật lý, quy trình mượn/trả, gia hạn, đặt chỗ (reservation), điều chuyển sách liên chi nhánh và hệ thống thông báo đa kênh.
    \item \textbf{Tối ưu hóa và Bảo mật:} Giải quyết các vấn đề bảo mật cơ bản như IDOR, SQL Injection (thông qua ORM), cấu hình CORS, thiết lập Rate Limiting và tối ưu hiệu suất truy vấn cơ sở dữ liệu (N+1 queries).
\end{itemize}

\subsection{Nhận xét và Đánh giá}
\textbf{Ưu điểm:}
\begin{itemize}
    \item Hệ thống có kiến trúc phân tán (Microservices) giúp tách biệt rõ ràng giữa logic nghiệp vụ thông thường (Backend) và xử lý tính toán nặng (AI Service), giúp dễ dàng mở rộng và bảo trì.
    \item Trải nghiệm người dùng được nâng cao rõ rệt thông qua các tính năng AI như tìm kiếm bằng ngôn ngữ tự nhiên và chatbot hỗ trợ 24/7.
    \item Hệ thống hoạt động ổn định, giao diện trực quan và phản hồi nhanh nhờ ứng dụng cơ chế caching và tối ưu hóa truy vấn.
\end{itemize}

\textbf{Hạn chế:}
\begin{itemize}
    \item Do triển khai trên hạ tầng miễn phí (Free Tier) của Vercel và Google Cloud Run, đôi khi hệ thống gặp tình trạng "Cold Start" (khởi động chậm) ở những yêu cầu đầu tiên hoặc bị giới hạn thời gian thực thi (timeout) với các tác vụ AI phức tạp.
    \item Thuật toán gợi ý cá nhân hóa chưa phát huy tối đa sức mạnh do lượng dữ liệu lịch sử mượn trả và đánh giá của người dùng trong môi trường thử nghiệm còn hạn chế (Cold-start problem cho Recommender System).
    \item Một số tính năng nâng cao như tóm tắt sách tự động đã được phát triển nhưng tạm thời vô hiệu hóa do chi phí API và giới hạn (Rate Limit) của mô hình LLM.
\end{itemize}

\subsection{Hướng phát triển}
Dựa trên những hạn chế và tiềm năng mở rộng, đồ án đề xuất các hướng phát triển trong tương lai:
\begin{itemize}
    \item \textbf{Nâng cấp hạ tầng triển khai:} Chuyển đổi sang hạ tầng trả phí (như AWS hoặc Google Cloud Platform chuyên dụng) để loại bỏ tình trạng Cold Start, tăng băng thông và hỗ trợ lượng người dùng lớn hơn.
    \item \textbf{Tối ưu hóa và mở rộng AI:} Nghiên cứu huấn luyện tinh chỉnh (Fine-tuning) các mô hình LLM mã nguồn mở nhỏ gọn (như Llama 3 hoặc Mistral) triển khai nội bộ để giảm phụ thuộc vào API bên thứ ba (Gemini), giúp tiết kiệm chi phí và tăng tốc độ phản hồi. Khôi phục và nâng cấp tính năng Function Calling cho Chatbot để hỗ trợ người dùng tra cứu thông tin cá nhân.
    \item \textbf{Hoàn thiện hệ thống gợi ý:} Tích hợp thêm các tín hiệu học máy nâng cao (chẳng hạn Deep Learning based recommender systems) và thu thập thêm dữ liệu hành vi người dùng (thời gian xem sách, click chuột) để gợi ý chính xác hơn.
    \item \textbf{Phát triển ứng dụng di động:} Xây dựng ứng dụng di động (Mobile App) đa nền tảng bằng React Native hoặc Flutter, tích hợp quét mã vạch/QR code để hỗ trợ thủ thư quản lý kho và độc giả thao tác nhanh chóng hơn.
\end{itemize}